% ============================================================
% 《数学分析基础》共享导言区
% ============================================================

\usepackage{amsmath,amssymb,amsthm}
\usepackage{geometry}
\geometry{a4paper, margin=2cm, top=2.5cm, bottom=2.5cm}

\usepackage[most]{tcolorbox}
\usepackage{graphicx, array, booktabs, multirow}
\usepackage{hyperref}
\hypersetup{colorlinks=true, linkcolor=blue!60!black, urlcolor=blue!60!black}
\usepackage{fancyhdr}
\usepackage{listings}
\usepackage{xcolor}
\usepackage{enumitem}
\usepackage{fontspec}

% ── 好题标记与汇总系统 ──
\usepackage[book]{goodproblems}

% ── 中文字体 ──
\setmainfont{Times New Roman}
\setCJKmainfont{PingFang SC}

% ── 图表路径 ──
\graphicspath{{figures/}}

% ── 代码块样式 ──
\lstdefinestyle{pystyle}{
  language=Python,
  basicstyle=\small\ttfamily,
  backgroundcolor=\color{gray!5},
  frame=single,
  numbers=left,
  numberstyle=\tiny\color{gray},
  breaklines=true,
  showstringspaces=false,
  keywordstyle=\color{blue},
  commentstyle=\color{green!40!black},
  stringstyle=\color{orange!70!black},
}
\lstset{style=pystyle}

% ── 章节标题 ──
\renewcommand{\thechapter}{\arabic{chapter}}

% ── 页眉页脚 ──
\pagestyle{fancy}
\fancyhf{}
\fancyhead[LE]{\leftmark}
\fancyhead[RO]{\rightmark}
\fancyfoot[C]{\thepage}

% ============================================================
% 定理类环境（颜色编码）
% ============================================================

% 定义框（绿色）
\newtcolorbox{definitionbox}[1][]{
  colback=green!5!white, colframe=green!60!black,
  fonttitle=\bfseries, title=定义~#1
}

% 定理框（蓝色）
\newtcolorbox{theorembox}[1][]{
  colback=blue!5!white, colframe=blue!60!black,
  fonttitle=\bfseries, title=定理~#1
}

% 命题框（紫色）
\newtcolorbox{propositionbox}[1][]{
  colback=purple!5!white, colframe=purple!60!black,
  fonttitle=\bfseries, title=命题~#1
}

% 例题框（橙色）
\newtcolorbox{examplebox}[1][]{
  colback=orange!5!white, colframe=orange!70!black,
  fonttitle=\bfseries, title=例题~#1
}

% 注意框（红色）
\newtcolorbox{notebox}[1][]{
  colback=red!5!white, colframe=red!60!black,
  fonttitle=\bfseries, title=注意~#1
}

% 进步宣言框（黄色）
\newtcolorbox{progressbox}{
  colback=yellow!10!white, colframe=yellow!60!orange,
  fonttitle=\bfseries\large, title=进步宣言
}

% 应用引导框（青色）
\newtcolorbox{applicationbox}{
  colback=teal!5!white, colframe=teal!60!black,
  fonttitle=\bfseries\large, title=学完本章你能做什么？
}

% 符号卡片（灰色）
\newenvironment{symbolcard}{%
  \par\vspace{6pt}
  \noindent\begin{tcolorbox}[colback=gray!3!white, colframe=gray!50!black, arc=2pt, title=📖 符号卡片]
  \renewcommand{\arraystretch}{1.3}
  \begin{tabular}{c|c|p{3.5cm}|p{5cm}}
  \toprule
  \bfseries 符号 & \bfseries LaTeX 写法 & \bfseries 读法 & \bfseries 含义 \\
  \midrule
}{%
  \bottomrule \end{tabular} \end{tcolorbox} \par\vspace{6pt}
}

% ============================================================
% 习题环境
% ============================================================

\newcounter{exercise}[chapter]
\renewcommand{\theexercise}{\thechapter.\arabic{exercise}}

\makeatletter
\newcommand{\difficulty}[1]{%
  \ifcase#1
    \textcolor{gray}{[送]}%
  \or
    \textcolor{gray}{[简]}%
  \or
    \textcolor{green!40!black}{[基]}%
  \or
    \textcolor{blue}{[普]}%
  \or
    \textcolor{orange}{[中]}%
  \or
    \textcolor{red}{[进]}%
  \or
    \textcolor{purple}{[拔]}%
  \or
    \textcolor{red!70!black}{[极]}%
  \or
    \textcolor{red!90!black}{[竞]}%
  \fi
}
\makeatother

\newenvironment{exercise}[1][0]{% 参数：难度等级 0-9
  \refstepcounter{exercise}
  \par\vspace{8pt}
  \noindent
  \textbf{练习 \theexercise}~\difficulty{#1}\quad
  \ignorespaces
}{\par\vspace{4pt}}

% 答案环境（放附录）
\newenvironment{answer}[1]{% 参数：练习编号
  \par\vspace{6pt}
  \noindent\textbf{#1}\quad
  \ignorespaces
}{\par\vspace{4pt}}

% ── 好题标记（最简单的用法）──
% 在 exercise 里面写 \好题 就行，不用起名字
\makeatletter
\newcommand{\好题}{%
  \edef\@gp@label{ex:\thechapter-\arabic{exercise}}%
  \label{prob:\@gp@label}%
  \GP@write{\string\GPentry{\theexercise}{\@gp@label}}%
}
\newcommand{\好答}{%
  \edef\@gp@label{ex:\thechapter-\arabic{exercise}}%
  \label{sol:\@gp@label}%
}
\makeatother

% ============================================================
% 其他设置
% ============================================================

% 图片不浮动（保持位置）
\setcounter{topnumber}{5}
\setcounter{bottomnumber}{5}
\setcounter{totalnumber}{10}

% 列表紧凑
\setlist{nosep}



% ── 自学教学设计环境 ──
\newtcolorbox{hintref}{colback=gray!5!white, colframe=gray!40!black,
  arc=2pt, left=4pt, right=4pt, top=3pt, bottom=3pt,
  fonttitle=\small\bfseries, title=提示}
\newtcolorbox{confidencebox}{colback=white, colframe=purple!40!black,
  arc=2pt, left=6pt, right=6pt, top=4pt, bottom=4pt,
  fonttitle=\small\bfseries, title=信心校准}
\newtcolorbox{retrievalgatebox}{colback=white, colframe=orange!60!black,
  arc=2pt, left=6pt, right=6pt, top=4pt, bottom=4pt,
  fonttitle=\small\bfseries, title=先想后看}
